%LTeX: language=fr-FR
\documentclass{article}
\usepackage{tasks}
\usepackage{amsmath,amssymb,amsfonts}
\usepackage[french]{babel}
\usepackage{enumitem}
\usepackage{currfile}
\usepackage{tikz}
\usetikzlibrary{arrows.meta}
\usepackage{tcolorbox}
\usepackage{xcolor}
\usepackage{fancyhdr}
\usepackage{array}
\usepackage[headsep=3mm,footskip=5mm,includefoot,includehead]{geometry}

\setlength{\parindent}{0cm}

\pagestyle{empty}
\input{\currfiledir ../def.tex}
\def \Document {Activité 1}


\geometry{left=1.5cm, right=1.5cm, top=1cm, bottom=1cm}
\setlength{\headheight}{14.0pt}

\newcolumntype{C}[1]{>{\centering\arraybackslash}m{#1}}

\newcommand{\rep}[1]{\textcolor{blue!70!black}{#1}}

% Courbe de la température au cours de la journée, avec un dessin en argument
\newcommand{\courbeT}[1]{%
  \begin{tikzpicture}[x=0.55cm, y=0.27cm]
    \draw[gray!25, very thin, xstep=1, ystep=1] (0,0) grid (24.5,16.5);
    \draw[gray!60, thin, xstep=2, ystep=2] (0,0) grid (24.5,16.5);
    \draw[-{Latex[length=2mm]}] (0,0) -- (25.3,0) node[right] {heure (h)};
    \draw[-{Latex[length=2mm]}] (0,0) -- (0,17.3) node[above] {température (°C)};
    \foreach \x in {2,4,...,24} \node[below] at (\x,0) {\small $\x$};
    \foreach \y in {2,4,...,16} \node[left] at (0,\y) {\small $\y$};
    \node[below left] at (0,0) {\small $0$};
    \draw[very thick] plot[smooth, tension=0.6] coordinates
      {(0,6) (3,4) (6,2) (8,6) (10,10) (12,13) (14,14) (16,13) (18,10) (21,8) (24,6)};
    #1
  \end{tikzpicture}%
}

\begin{document}
\pagestyle{fancy}
\fancyhf{}
\fancyhead[L]{\textbf{\Niveau}}
\fancyhead[C]{\textbf{\ChapitreCourt}}
\fancyhead[R]{\textbf{\Document\ -- Correction}}

\normalsize

\begin{center}
  {\Large\bfseries Correction : fonctions du quotidien}
\end{center}

% =====================================================
\subsection*{Partie A : la température au cours d'une journée}

Une station météo a enregistré la température à Lyon le 2 novembre, de 0 h à 24 h.

\begin{center}
  \courbeT{
    \draw[dashed, red!70!black] (0,10) -- (18,10) (10,0) -- (10,10) (18,0) -- (18,10);
    \draw[line width=2.5pt, red!70!black] (10,0) -- (18,0);
    \draw[dashed, blue!70!black] (0,6) -- (24,6) (8,0) -- (8,6);
  }
\end{center}

\begin{enumerate}
  \item Quelle était la température à 6 h~? À 14 h~? À 18 h~?

  \rep{À 6 h : $2$~°C. \quad À 14 h : $14$~°C. \quad À 18 h : $10$~°C.}

  \item À quelle(s) heure(s) faisait-il $10$~°C~? Et $6$~°C~?

  \rep{$10$~°C à 10 h et à 18 h. \quad $6$~°C à 0 h, à 8 h et à 24 h.}

  \item A-t-il fait $16$~°C ce jour-là~? \rep{Non : la température maximale est $14$~°C.}

  \item Peut-on lire sur ce graphique la température à 26 h~? Pourquoi~? \rep{Non : le graphique ne va que de 0 h à 24 h.}

  \item Pendant quelle période de la journée la température était-elle supérieure ou égale à $10$~°C~? \rep{De 10 h à 18 h, soit $x \in [10\,;\,18]$.}
\end{enumerate}

\begin{tcolorbox}[colback=white, colframe=black, title=Vocabulaire]
  À chaque heure $x$ de l'intervalle $[0\,;\,24]$, la courbe associe \textbf{une seule} température, notée $T(x)$ : on dit que $T$ est une \textbf{fonction}. L'intervalle $[0\,;\,24]$ est son \textbf{ensemble de définition}.

  Comme $T(6) = 2$, on dit que $2$ est \textbf{l'image} de $6$ par $T$, et que $6$ est \textbf{un antécédent} de $2$ par $T$.
\end{tcolorbox}

\begin{enumerate}[resume]
  \item Compléter : \quad $T(14) = \rep{14}$ \qquad $T(\rep{6}) = 2$ \qquad Les antécédents de $6$ par $T$ sont \rep{$0$, $8$ et $24$}.
\end{enumerate}

% =====================================================
\subsection*{Partie B : l'abonnement au club de sport}

Un club de sport fait payer une carte de membre à $10$~€, puis $4$~€ par séance. On note $f(x)$ le prix payé, en euros, pour $x$ séances.

\begin{enumerate}[resume]
  \item Compléter le tableau de valeurs.

  \begin{center}
    \renewcommand{\arraystretch}{1.3}
    \begin{tabular}{|C{4.2cm}|*{6}{C{1.4cm}|}}
      \hline
      Nombre de séances $x$ & $0$ & $1$ & $2$ & $3$ & $4$ & $5$ \\ \hline
      Prix $f(x)$ (en €) & \rep{$10$} & \rep{$14$} & \rep{$18$} & \rep{$22$} & \rep{$26$} & \rep{$30$} \\ \hline
    \end{tabular}
  \end{center}

  \item Exprimer $f(x)$ en fonction de $x$. \rep{$f(x) = 4x + 10$}

  \item Calculer $f(8)$. Que représente ce nombre~? \rep{$f(8) = 4 \times 8 + 10 = 42$ : le prix de 8 séances est $42$~€.}

  \item Léa a payé $50$~€. Combien de séances a-t-elle faites~? Quel est donc l'antécédent de $50$ par $f$~?

  \rep{$4x + 10 = 50 \iff 4x = 40 \iff x = 10$. Léa a fait $10$ séances : l'antécédent de $50$ par $f$ est $10$.}
\end{enumerate}

\end{document}
